% Chapter — Solar radiation pressure and conical shadow (FR-7). Follows the
% mandatory FR-29 six-section template: purpose; assumptions and domain of
% validity; derivation; discretization and implementation; validation
% evidence; references. The equation labels eq:srp:pressure,
% eq:srp:cannonball, eq:srp:appgeom, eq:srp:overlap, eq:srp:nu, and
% eq:srp:boundaries are echoed verbatim in comments in
% cpp/src/models/srp.cpp; eq:srp:umbracone and eq:srp:penumbracone are
% echoed in the independent analytic reference of
% cpp/tests/test_srp.cpp (FR-29 equation-label-to-code traceability);
% renaming them breaks that trace.
\chapter{Solar Radiation Pressure and Conical Shadow}
\label{ch:srp}

\section{Purpose}
\label{sec:srp:purpose}

This chapter documents the solar-radiation-pressure model (FR-7): the
attitude-independent cannonball acceleration with ephemeris Sun distance,
scaled by an illumination fraction $\nu \in [0, 1]$ from a dual-cone
(umbra/penumbra) conical shadow model that any configured spherical body
--- including the current central body --- may cast. The shadow function
is a separately callable, separately testable component: besides scaling
the SRP force it supplies the signed boundary functions that the
event-detection framework of Chapter~\ref{ch:events} screens for
eclipse-entry and -exit events (FR-12). The model is implemented in
\texttt{cpp/src/models/srp.cpp} with its interface in
\texttt{cpp/include/star/models/srp.hpp}; Sun and occulter positions come
from the ephemeris evaluator of Chapter~\ref{ch:ephemeris}, and all body
radii and the spacecraft ballistic parameter $C_R A / m$ are supplied by
the caller.

\section{Assumptions and domain of validity}
\label{sec:srp:domain}

Assumptions:
\begin{enumerate}
  \item \textbf{Cannonball spacecraft.} The SRP force is
    attitude-independent, characterized by the single lumped parameter
    $C_R A / m$: an effective cross-section $A$, mass $m$, and radiation
    pressure coefficient $C_R$ that absorbs the absorbed/reflected
    momentum accounting ($C_R = 1$ for a perfect absorber, up to $2$ for
    a flat specular reflector facing the Sun; a sphere with diffuse
    reflectivity $\varepsilon$ has $C_R = 1 + \tfrac{4}{9}\varepsilon$,
    Montenbruck and Gill~\cite{montenbruck2000}, Sec.~3.4). The N-plate
    upgrade path of FR-7 replaces this parameterization without touching
    the shadow model.
  \item \textbf{Inverse-square flux from the nominal irradiance.} The
    solar flux is $\Phi(d) = \Phi_{1\,\mathrm{au}} (\mathrm{au}/d)^{2}$
    with the IAU 2015 Resolution B3 nominal total solar irradiance
    $\Phi_{1\,\mathrm{au}} = \qty{1361}{\watt\per\metre\squared}$ (Pr{\v s}a
    et al.~\cite{prsa2016}). The measured solar-cycle TSI variation is
    $\sim \qty{1.6}{\watt\per\metre\squared}$ peak-to-peak
    ($\sim 0.12\%$) about a minimum of
    $\qty{1360.8 \pm 0.5}{\watt\per\metre\squared}$ (Kopp and
    Lean~\cite{kopp2011}); using the constant nominal bounds the SRP
    force error at the same relative level, far below the uncertainty of
    any configured $C_R$.
  \item \textbf{Geometric Sun position.} As in
    Chapter~\ref{ch:thirdbody}, the Sun position is geometric (FR-23
    light-time policy). Light-time and aberration displace the apparent
    Sun by $\sim 10^{-4}$ of its distance, a direction error of
    $\sim \qty{e-4}{\radian}$ and a flux error of $\sim 10^{-4}$
    relative --- both absorbed by the $C_R$ uncertainty.
  \item \textbf{Spherical, opaque, airless occulters.} Each occulting
    body is a sphere of caller-supplied radius. Neglected against the
    Earth case: oblateness (the $\sim 1/298$ flattening perturbs the
    shadow-boundary radius by up to $\sim \qty{21}{\kilo\metre}$,
    shifting LEO crossing times by up to a few seconds depending on
    crossing geometry) and atmospheric refraction/extinction (which
    softens the real penumbra by an altitude-dependent amount). These
    are physical-fidelity bounds of the model, not of its
    implementation; the Phase~3 criterion gates the implementation
    against the model's own exact geometry.
  \item \textbf{Uniform solar disk.} Solar limb darkening is neglected:
    the penumbra illumination fraction is the uncovered \emph{area}
    fraction of a uniformly bright disk (the standard model,
    Montenbruck and Gill~\cite{montenbruck2000}, Sec.~3.4).
  \item \textbf{Flat-disk overlap in the penumbra interior.} The
    umbra/penumbra/annular \emph{boundaries} follow from exact sphere
    tangency (Section~\ref{sec:srp:derivation}); only the partial-phase
    \emph{interior} value of $\nu$ treats the two apparent disks as
    planar circles of radii equal to their angular radii. For a LEO
    spacecraft eclipsed by the Earth the occulter's angular radius
    ($\sim \qty{1.2}{\radian}$) is large, but the solar disk
    ($\sim \qty{4.7e-3}{\radian}$) sees the occulter limb as locally
    straight to $\sim a_{s}/\tan a_{o} \lesssim 0.3\%$; the induced
    $\nu$ error is below $10^{-2}$ of the penumbra transition, which
    lasts seconds and carries $\lesssim 10^{-7}\,\mathrm{m/s^2}$ of
    acceleration --- negligible against assumption 4.
\end{enumerate}

Domain of validity: spacecraft outside both spheres
($\norm{\vv{r}_{\odot} - \vv{r}} > R_{\odot}$,
$\norm{\vv{r}_{o} - \vv{r}} > R_{o}$), any occulter/Sun geometry
including the annular case ($a_{o} < a_{s}$), occulters at any distance.
Multiple occulting bodies are composed by the force-assembly layer
(evaluating $\nu$ per occulter); this chapter defines the single-occulter
function.

Out-of-domain behavior, matching the code exactly: if the spacecraft is
inside an occulting sphere, the $\arcsin$ argument is clamped to 1
($a_{o} = \pi/2$) and the function still returns a well-defined value
continuous with the exterior; sub-surface states are excluded by
ground-impact eventing (FR-12), so no throw is warranted in the hot
loop. If the spacecraft coincides with the Sun's center the direction
$\hat{\vv{s}}$ is undefined and IEEE NaNs propagate. $\nu$ is exactly
$1$ in full sunlight and exactly $0$ in total umbra by construction
(returned as literal constants, not computed expressions).

\section{Derivation}
\label{sec:srp:derivation}

\subsection{Radiation pressure and the cannonball acceleration}

A photon flux of energy-flux density $\Phi$ carries momentum-flux
density $\Phi / c$ (each photon of energy $E$ carries momentum $E/c$),
so a fully absorbing surface normal to the flux experiences pressure
$P = \Phi / c$. With the inverse-square flux of assumption 2, the
radiation pressure at heliocentric distance $d$ is
\begin{equation}
  P(d) = \frac{\Phi_{1\,\mathrm{au}}}{c}
         \left( \frac{\mathrm{au}}{d} \right)^{2}
       = P_{1\,\mathrm{au}} \left( \frac{\mathrm{au}}{d} \right)^{2},
  \qquad
  P_{1\,\mathrm{au}} = \frac{\qty{1361}{\watt\per\metre\squared}}
                            {\qty{299792458}{\metre\per\second}}
                     \approx \qty{4.5398e-6}{\newton\per\metre\squared},
  \label{eq:srp:pressure}
\end{equation}
with $c$ exact by the SI definition of the metre~\cite{bipm2019} and the
astronomical unit exact per IAU 2012 Resolution B2~\cite{iau2012b2}. The
cannonball force is $P C_R A$ directed along the Sun-to-spacecraft unit
vector, giving the acceleration
\begin{equation}
  \boxed{\;
  \vv{a}_{\mathrm{srp}}
    = \nu \, P_{1\,\mathrm{au}}
      \left( \frac{\mathrm{au}}{d} \right)^{2}
      \frac{C_R A}{m} \, \hat{\vv{s}},
  \qquad
  \hat{\vv{s}} = \frac{\vv{r} - \vv{r}_{\odot}}
                      {\norm{\vv{r} - \vv{r}_{\odot}}},
  \quad
  d = \norm{\vv{r} - \vv{r}_{\odot}},
  \;}
  \label{eq:srp:cannonball}
\end{equation}
where $\vv{r}$ and $\vv{r}_{\odot}$ are the spacecraft and Sun positions
in a common frame about a common origin and $\nu$ is the illumination
fraction derived below. \textbf{Direction convention, stated
explicitly:} $\hat{\vv{s}}$ points \emph{from the Sun to the
spacecraft} --- the propagation direction of the photons --- so with
$C_R > 0$ the acceleration pushes the spacecraft away from the Sun. The
model follows Montenbruck and Gill~\cite{montenbruck2000}, Sec.~3.4.

\subsection{Apparent-disk shadow geometry}

Seen from the spacecraft, the Sun and the occulting body are disks with
apparent angular radii and angular center separation
\begin{equation}
  a_{s} = \arcsin \frac{R_{\odot}}{\norm{\vv{r}_{\odot} - \vv{r}}},
  \qquad
  a_{o} = \arcsin \frac{R_{o}}{\norm{\vv{r}_{o} - \vv{r}}},
  \qquad
  \theta = \angle\bigl( \vv{r}_{\odot} - \vv{r},\;
                        \vv{r}_{o} - \vv{r} \bigr),
  \label{eq:srp:appgeom}
\end{equation}
with $R_{\odot}$ the solar radius, $R_{o}$ the occulter radius, and
$\vv{r}_{o}$ the occulter position. These angular quantities are exact
for spheres. Four regimes follow immediately from disk containment and
tangency:
\begin{itemize}
  \item $\theta \geq a_{s} + a_{o}$: the disks are disjoint --- full
    sunlight, $\nu = 1$.
  \item $\theta \leq a_{o} - a_{s}$ (possible only when
    $a_{o} \geq a_{s}$): the solar disk lies entirely inside the
    occulter's --- total umbra, $\nu = 0$.
  \item $\theta \leq a_{s} - a_{o}$ (possible only when
    $a_{s} > a_{o}$): the occulter's disk lies entirely inside the
    Sun's --- annular eclipse (antumbra).
  \item otherwise: partial overlap --- penumbra.
\end{itemize}

For the annular case the covered area is the full occulter disk, so on
the uniform-disk assumption
\begin{equation}
  \nu_{\mathrm{annular}} = 1 - \left( \frac{a_{o}}{a_{s}} \right)^{2},
  \label{eq:srp:annular}
\end{equation}
independent of $\theta$ while the containment lasts.

For the partial case, the covered area is the lens common to two planar
disks of radii $a_{s}$, $a_{o}$ with center separation $\theta$
(assumption 6). Place the Sun-disk center at the origin and the
occulter-disk center at distance $\theta$ on the axis. The two circles
$x^{2} + y^{2} = a_{s}^{2}$ and $(x - \theta)^{2} + y^{2} = a_{o}^{2}$
intersect on the chord at abscissa and half-height
\begin{equation}
  x_{c} = \frac{\theta^{2} + a_{s}^{2} - a_{o}^{2}}{2\theta},
  \qquad
  y_{c} = \sqrt{a_{s}^{2} - x_{c}^{2}}
  \label{eq:srp:chord}
\end{equation}
(subtract the circle equations to eliminate $y^{2}$ and solve for $x$).
The lens is the union of two circular segments, one from each disk, cut
off by the common chord. A segment of a disk of radius $\rho$ whose
chord lies at center distance $h$ has area $\rho^{2} \arccos(h/\rho) -
h \sqrt{\rho^{2} - h^{2}}$ (the circular sector minus the two triangles
between the radii and the chord). With $h = x_{c}$ for the solar disk
and $h = \theta - x_{c}$ for the occulter disk, and
$\sqrt{\rho^{2} - h^{2}} = y_{c}$ on the shared chord in both cases, the
overlap area is
\begin{equation}
  A_{\cap} = a_{s}^{2} \arccos \frac{x_{c}}{a_{s}}
           + a_{o}^{2} \arccos \frac{\theta - x_{c}}{a_{o}}
           - \theta\, y_{c},
  \label{eq:srp:overlap}
\end{equation}
where the two triangle contributions combined to $\theta\, y_{c}$
($x_{c} y_{c} + (\theta - x_{c}) y_{c}$). This is the standard
occultation formula (Montenbruck and Gill~\cite{montenbruck2000},
Sec.~3.4); the construction above is reproduced so no step is external
to this document. Note $x_{c}$ may be negative (chord beyond the solar
disk center, $\arccos$ obtuse) --- the formula remains valid on
$[-1, 1]$ arguments. Collecting all regimes, the illumination fraction
is
\begin{equation}
  \boxed{\;
  \nu(\theta, a_{s}, a_{o}) =
  \begin{cases}
    1, & \theta \geq a_{s} + a_{o}, \\[2pt]
    0, & \theta \leq a_{o} - a_{s}, \\[2pt]
    1 - \left( a_{o} / a_{s} \right)^{2},
       & \theta \leq a_{s} - a_{o}, \\[2pt]
    1 - A_{\cap} / \bigl( \pi a_{s}^{2} \bigr), & \text{otherwise.}
  \end{cases}
  \;}
  \label{eq:srp:nu}
\end{equation}
$\nu$ is continuous across every regime boundary: at
$\theta = a_{s} + a_{o}$, $x_{c} = a_{s}$ and both $\arccos$ terms and
$y_{c}$ vanish, so $A_{\cap} = 0$ and $\nu \to 1$; at
$\theta = a_{o} - a_{s}$, $x_{c} = -a_{s}$, $y_{c} = 0$,
$A_{\cap} = \pi a_{s}^{2}$ and $\nu \to 0$; and at
$\theta = a_{s} - a_{o}$, $x_{c} = a_{s}$ and
$\theta - x_{c} = -a_{o}$, so $y_{c} = 0$,
$A_{\cap} = \pi a_{o}^{2}$, and $\nu$ matches
equation~\eqref{eq:srp:annular}.

\subsection{Signed boundary functions for event detection}

The event framework (Chapter~\ref{ch:events}) screens scalar functions
for sign changes, so the two shadow boundaries are exposed directly:
\begin{equation}
  g_{u} = \theta - (a_{o} - a_{s}),
  \qquad
  g_{p} = \theta - (a_{s} + a_{o}).
  \label{eq:srp:boundaries}
\end{equation}
$g_{u} < 0$ exactly in the total umbra ($g_{u} > 0$ everywhere when the
geometry is annular, where no umbra exists), and $g_{p} < 0$ exactly
where any part of the Sun is occulted ($\nu < 1$). Both are smooth in
the spacecraft position through each boundary and cross zero
transversally along an orbit, unlike $\nu$ itself, which is constant
outside the penumbra and therefore useless for sign screening.

\subsection{The tangent cones and the analytic timing reference}
\label{sec:srp:cones}

The disk-tangency boundaries of equation~\eqref{eq:srp:nu} coincide
exactly, for spherical bodies, with the classical umbra and penumbra
cones --- the construction the validation test uses as an independent
reference. Let the occulter center sit at the origin and the Sun center
at distance $D$ on the axis, with $R_{\odot} > R_{o}$.

\textbf{Umbra cone.} The common \emph{external} tangent lines touch
both spheres on the same side of the axis and intersect the axis at the
apex $V_{u}$ beyond the occulter, at distance $\ell_{u}$ from the
occulter center. Similar right triangles from the apex to each tangent
point (legs $R_{o}$ and $R_{\odot}$, hypotenuses along the same line)
give $\sin\alpha_{u} = R_{o}/\ell_{u} = R_{\odot}/(D + \ell_{u})$, so
\begin{equation}
  \ell_{u} = \frac{R_{o} D}{R_{\odot} - R_{o}},
  \qquad
  \sin\alpha_{u} = \frac{R_{\odot} - R_{o}}{D},
  \label{eq:srp:umbracone}
\end{equation}
where $\alpha_{u}$ is the cone half-angle. A point $P$ between the
occulter and the apex is on the cone surface exactly when the
tangent line through it grazes both spheres; looking along that line
toward the Sun, the ray makes angle $a_{o}$ with the occulter-center
direction and angle $a_{s}$ with the Sun-center direction, \emph{with
both centers on the same side of the ray}, so the center separation is
$\theta = a_{o} - a_{s}$ --- precisely the internal disk tangency
$g_{u} = 0$. The umbra-cone surface and the $\nu = 0$ boundary are the
same set.

\textbf{Penumbra cone.} The common \emph{internal} tangent lines cross
the axis between the bodies at the apex $V_{p}$, at distance $\ell_{p}$
from the occulter center toward the Sun, and graze the two spheres on
opposite sides of the axis: $\sin\alpha_{p} = R_{o}/\ell_{p} =
R_{\odot}/(D - \ell_{p})$, so
\begin{equation}
  \ell_{p} = \frac{R_{o} D}{R_{\odot} + R_{o}},
  \qquad
  \sin\alpha_{p} = \frac{R_{\odot} + R_{o}}{D}.
  \label{eq:srp:penumbracone}
\end{equation}
From a point on this cone (behind the occulter), the grazing ray toward
the Sun has the two centers on \emph{opposite} sides, so
$\theta = a_{o} + a_{s}$ --- the external disk tangency $g_{p} = 0$.

For a circular orbit of radius $a$ about the occulter with the Sun
direction fixed along $+\hat{\vv{x}}$ and the orbit in a plane
containing the Sun line, the position is $(a\cos\psi, a\sin\psi)$ and
the umbra-cone crossing condition (distance from the axis equals the
cone radius at that axial station, on the shadow side) reads
$\lvert a \sin\psi \rvert = \tan\alpha_{u}\,(a\cos\psi + \ell_{u})$.
Squaring and substituting $u = \cos\psi$ yields the quadratic
\begin{equation}
  a^{2} (1 + \tan^{2}\alpha_{u})\, u^{2}
  + 2 a\, \ell_{u} \tan^{2}\alpha_{u}\, u
  + \bigl( \ell_{u}^{2} \tan^{2}\alpha_{u} - a^{2} \bigr) = 0,
  \label{eq:srp:conequadratic}
\end{equation}
whose two roots have a negative product whenever
$\ell_{u} \tan\alpha_{u} < a$ (the cone still overlaps the orbit
radius); the negative root is the shadow-side crossing
($\cos\psi < 0$), and the squaring artifact (the mirror crossing on the
sunlit side) is the positive root. Entry and exit are then
$\psi = \arccos u$ and $2\pi - \arccos u$, converted to time by the
orbit rate. The penumbra times follow identically with
$\ell_{p} \to$ apex on the $+\hat{\vv{x}}$ side, i.e.\
$\lvert a \sin\psi \rvert = \tan\alpha_{p}\,(\ell_{p} - a\cos\psi)$.
This closed form shares no code and no formulation with
equations~\eqref{eq:srp:appgeom}--\eqref{eq:srp:boundaries} beyond the
sphere geometry itself, which is what makes the criterion-3 timing
comparison meaningful.

\section{Discretization and implementation}
\label{sec:srp:implementation}

There is no discretization: both functions are pointwise algebraic maps
evaluated at integrator-requested states. Implementation notes:
\begin{enumerate}
  \item \textbf{Module and traceability.} The model lives in
    \texttt{cpp/src/models/srp.cpp}; source comments echo
    \texttt{eq:srp:pressure}, \texttt{eq:srp:cannonball},
    \texttt{eq:srp:appgeom}, \texttt{eq:srp:overlap},
    \texttt{eq:srp:nu}, and \texttt{eq:srp:boundaries} verbatim. The
    analytic cone reference (equations~\eqref{eq:srp:umbracone},
    \eqref{eq:srp:penumbracone}, \eqref{eq:srp:conequadratic}) is
    implemented \emph{only} in \texttt{cpp/tests/test\_srp.cpp}, echoed
    there, so the validation reference cannot share a defect with the
    model.
  \item \textbf{Shared apparent geometry.} One internal helper computes
    $(a_{s}, a_{o}, \theta)$ for the fraction and both boundary
    functions, so the three public functions cannot disagree about the
    geometry.
  \item \textbf{Angle numerics.} $\theta$ is evaluated as
    $\operatorname{atan2}(\norm{\vv{u} \times \vv{v}},
    \vv{u} \cdot \vv{v})$, which is accurate for all separations,
    unlike $\arccos$ of the normalized dot product, which loses
    precision near $0$ and $\pi$ --- and eclipse geometry lives near
    $\theta \approx a_{o}$, which for LEO is $\mathcal{O}(1)$ but must
    resolve $\mathcal{O}(a_{s})$ differences. $\arcsin$ arguments are
    clamped to $[0, 1]$ (out-of-domain rule above); $\arccos$ arguments
    in equation~\eqref{eq:srp:overlap} are clamped to $[-1, 1]$ and the
    chord height uses $\sqrt{\max(0, \cdot)}$, protecting
    boundary-adjacent states from sub-ulp excursions outside the real
    domain.
  \item \textbf{Exact 0 and 1.} The full-sunlight and total-umbra
    branches return the literals $1.0$ and $0.0$ before any area
    arithmetic runs, so $\nu$ is bit-exact in both regimes ---
    the property the determinism gates and the force-by-source log
    channels rely on. The umbra case of the acceleration is exactly
    zero because $\nu = 0$ multiplies the finite coefficient.
  \item \textbf{Conditioning of the partial branch.} In
    equation~\eqref{eq:srp:overlap} the occulter-disk term
    $a_{o}^{2} \arccos\bigl((\theta - x_{c})/a_{o}\bigr)$ evaluates its
    $\arccos$ near argument $1$ when the occulter disk is much larger
    than the solar disk (LEO Earth shadow: $a_{o} \approx
    \qty{1.2}{\radian}$ vs.\ $a_{s} \approx \qty{4.7e-3}{\radian}$),
    amplifying ulp-level angle errors by $a_{o}^{2} / \bigl( \pi
    a_{s}^{2} \sqrt{1 - ((\theta - x_{c})/a_{o})^{2}} \bigr) \sim
    10^{6}$--$10^{7}$ after the $\pi a_{s}^{2}$ normalization. The
    double-precision partial-phase $\nu$ therefore carries an absolute
    error of order $10^{-9}$ at such geometries (observed worst
    \num{1.0e-9} at the committed penumbra states), while small-occulter
    partial and annular geometries stay at $10^{-12}$. This is retained
    deliberately: the formula is the cited standard, the error is
    $\sim 10^{-9}$ of the SRP force during a seconds-long transit ---
    far below the $C_R$ knowledge of assumption 1 and the physical
    fidelity bounds of assumptions 4--6 --- and the regime boundaries
    and exact $0$/$1$ branches are unaffected. The golden tolerance
    (\num{1e-8} absolute on partial values) encodes exactly this
    bound.
  \item \textbf{Constants.} $c$, au, $\Phi_{1\,\mathrm{au}}$,
    $P_{1\,\mathrm{au}} = \Phi_{1\,\mathrm{au}}/c$ (a compile-time
    quotient of the two cited constants), and $R_{\odot}$ live in
    \texttt{cpp/include/star/constants.hpp}, each cited at its
    definition. Occulter radii are caller-supplied (test geometry uses
    representative Earth/Moon values; physical values are wired by the
    force-composition layer with the rest of the body data).
  \item \textbf{Determinism.} No per-call allocation; fixed evaluation
    order. The libm calls ($\arcsin$, $\arccos$,
    $\operatorname{atan2}$) are the only non-correctly-rounded
    operations and thus the cross-platform divergence surface (D-10);
    the golden tolerances (\texttt{tests/golden/srp/manifest.toml})
    budget for their few-ulp spread.
\end{enumerate}

\section{Validation evidence}
\label{sec:srp:validation}

Table~\ref{tab:srp:validation} lists the tests that constitute this
chapter's validation evidence. Each test identifier is machine-checked
to exist in the test tree by \texttt{scripts/lint\_chapters.py}; golden
provenance and tolerances are recorded in
\texttt{tests/golden/srp/manifest.toml}.

\begin{table}[htbp]
  \centering
  \caption{Validation evidence for the SRP and shadow model (doctest,
    \texttt{cpp/tests/test\_srp.cpp}).}
  \label{tab:srp:validation}
  \begin{tabular}{lp{8.6cm}}
    \toprule
    Test ID & Acceptance basis \\
    \midrule
    \testid{srp_shadow_cone_timing} &
      Phase~3 exit criterion 3: umbra entry and exit times of a
      circular-LEO shadow crossing, root-found on
      equation~\eqref{eq:srp:boundaries} by bisection along the analytic
      orbit, match the closed-form cone times of
      equation~\eqref{eq:srp:conequadratic} to
      $< \qty{0.1}{\second}$; penumbra times likewise; and the conical
      entry/exit differ from a cylindrical shadow's by
      $> \qty{1}{\second}$ (expected $\approx \qty{4}{\second}$),
      guarding against silently shipping the cylinder. \\
    \testid{srp_shadow_fraction_golden} &
      Equation~\eqref{eq:srp:nu} vs.\ extended-precision references at
      11 committed geometries covering every branch: exact $1$/$0$ in
      full sun and umbra (bit equality), penumbra depths, annular and
      off-axis annular, partial lunar occultation, Moon-occulter umbra
      (abs $10^{-8}$ on partial values per the conditioning bound of
      Section~\ref{sec:srp:implementation}; observed worst
      \num{1.0e-9}). \\
    \testid{srp_cannonball_accel_golden} &
      Equation~\eqref{eq:srp:cannonball} vs.\ extended-precision
      references at 5 committed states, norm-relative $10^{-12}$,
      including the exact-zero umbra case and an inverse-square check
      at Mars heliocentric distance. \\
    \testid{srp_shadow_fraction_continuity} &
      Property test: along the reference shadow crossing, $\nu$ is $1$
      before the penumbra, $0$ inside the umbra, strictly between $0$
      and $1$ at the penumbra midpoint, within $[0,1]$ everywhere, and
      changes by $< 10^{-3}$ between samples \qty{1}{\milli\second}
      apart (continuity; the physical penumbra transit is
      $\approx \qty{8}{\second}$). \\
    \bottomrule
  \end{tabular}
\end{table}

\section{References}
\label{sec:srp:references}

This chapter relies on Montenbruck and Gill~\cite{montenbruck2000},
Sec.~3.4, for the cannonball SRP model and the apparent-disk occultation
formulation; on Pr{\v s}a et al.~\cite{prsa2016} for the IAU 2015
Resolution B3 nominal total solar irradiance and nominal solar radius;
on Kopp and Lean~\cite{kopp2011} for the measured TSI value and its
variability bound; on IAU 2012 Resolution B2~\cite{iau2012b2} for the
exact astronomical unit; and on the BIPM SI brochure~\cite{bipm2019}
for the exact speed of light. Full bibliographic details appear in the
bibliography.
